\section{Galaxy Clusters}

Galaxy clusters are the largest gravitationally bound structures in the Universe, 
comprising hundreds to thousands of galaxies, 
vast amounts of hot intracluster gas that emits X-rays, and of dark matter.
Typically, galaxy clusters span between $1$ and $5$~Mpc in diameter, 
in comparison individual galaxies (such as the Milky way), are on the order of $10$--$100$~kpc. Their total 
masses are in the range $10^{14}$ to $10^{15}~M_\odot$. 

The composition of a galaxy cluster includes three main components. First, the 
galaxies themselves, which contribute only a small fraction of the total mass. 
Second, the hot intracluster gas (plasma), which dominates the baryonic mass and is heated to temperatures 
of $10^7$--$10^8$~K, emits strongly in X-rays via thermal bremsstrahlung. Finally, dark matter, 
which is about about 80\% of the total mass, a mysterious form of matter that has so far escaped direct detection with any type of 
telescope, but makes its presence felt through its gravitational pull on the galaxies and hot gas.

Clusters cannot be made only of galaxies, as the observed gravitational potential 
(measured via the velocities of galaxies and the temperature distribution of the gas) 
far exceeds the mass inferred from visible matter alone. This discrepancy highlights the 
crucial role of both the hot intracluster medium (ICM) and dark matter in cluster dynamics 
and structure.

In different observable wavelengths, galaxy clusters exhibit different features.
In optical light, they appear as dense concentrations of galaxies. In contrast, X-ray observations 
reveal emission from the hot ICM, outlining the gravitational potential of the cluster. 
Observatories such as Chandra, XMM-Newton, and eROSITA are sensitive to this emission in the $0.1$--$10$~keV energy range.

The X-ray emission from the ICM is primarily due to thermal bremsstrahlung (free-free emission), which depends on the density and temperature of the gas. The spatial distribution of the X-ray emitting gas is often modeled using the $\beta$-model (see \cite{Sarazin2002_NED}):

\[
\rho(r) = \rho_0 \left[ 1 + \left( \frac{r}{r_c} \right)^2 \right]^{-3\beta/2},
\]

where $\rho_0$ is the central gas density, $r_c$ is the core radius, and $\beta$ is a slope parameter. By projecting this profile along the line of sight, one obtains the X-ray surface brightness profile, which can be fit to observational data.

Under the assumption of an isothermal gas in hydrostatic equilibrium, the total gas mass $M_g$ of a galaxy cluster can be estimated using the $\beta$-model as:

\[
M_g = \pi^{3/2} \, \rho_{g0} \, r_c^3 \, \frac{\Gamma\left( \frac{3\beta - 1}{2} \right)}{\Gamma\left( \frac{3\beta}{2} \right)} \quad \text{for } \beta > 1
\]

Alternatively, in terms of the central gas number density $n_0$,

\[
M_g = 3.15 \times 10^{12} \, M_\odot \left( \frac{n_0}{10^{-3} \, \text{cm}^{-3}} \right) \left( \frac{r_c}{0.25 \, \text{Mpc}} \right)^3 \frac{\Gamma\left( \frac{3\beta - 1}{2} \right)}{\Gamma\left( \frac{3\beta}{2} \right)}
\]

where:
\begin{itemize}
    \item $n_0$ is the central gas number density,
    \item $r_c$ is the core radius of the cluster,
    \item $\beta$ is the slope parameter from the $\beta$-model,
    \item $\Gamma$ is the gamma function,
    \item $\rho_{g0}$ is the central gas density.
\end{itemize}

These relations allow us to estimate the 
gas and total mass profiles of a cluster from observed X-ray 
surface brightness and gas temperature, providing key information 
for understanding the gravitational mass distribution in galaxy 
clusters.


\section{Convolutional Neural Networks (CNNs) for Image Analysis}

Convolutional Neural Networks (CNNs) are a class of deep learning 
models particularly effective for image analysis. In this lab, due to the nature of our data (iamages of galaxy clusters)
we will build and train CNNs.
For more theoretical background, 
see Goodfellow et al.~\cite{Goodfellow-et-al-2016} and the TensorFlow 
CNN tutorial\footnote{\url{https://www.tensorflow.org/tutorials/images/cnn}}. 


A standard CNN typically consists of the following components:
\begin{itemize}
    \item \textbf{Input layer:} Accepts the input image data.
    \item \textbf{Convolutional layers (Conv2D):} Apply filters (kernels) to 
    extract spatial features from the input.
    \item \textbf{Activation functions:} Introduce nonlinearity (commonly ReLU).
    \item \textbf{Pooling layers:} Downsample feature maps, typically via max pooling.
    \item \textbf{Dropout layers:} Reduce overfitting by randomly deactivating 
    neurons during training.
    \item \textbf{Fully connected layers:} Integrate features to make final predictions.
    \item \textbf{Output layer:} Produces the classification result.
\end{itemize}

The core of a CNN is the \textbf{convolutional layer}, this layer applies a set of 
learnable filters (kernels) on the input. Each filter slides across the 
input image (or previous feature map), performing an element-wise 
multiplication and summing the result to produce a \emph{feature map}. 
Important hyperparameters include:

\begin{itemize}
    \item \textbf{Kernel size:} The spatial dimensions of the filter (e.g., $3 \times 3$).
    \item \textbf{Filters:} The number of different filters (feature maps) in the layer.
    \item \textbf{Stride:} The step size the filter moves with; higher stride means more downsampling.
    \item \textbf{Padding:} Controls if the input is padded (e.g., with zeros) at the borders.
\end{itemize}

A key difference between CNNs and fully connected (dense) neural networks is that
CNNs exploit \emph{local spatial correlations} in image data through weight sharing.
Weight sharing in CNNs means that the same set of filter weights is 
applied across the entire input image as the filter moves. 
This enables the network 
to detect the same pattern regardless of its position in the image.
Whereas, on the other hand, dense networks connect every input to every 
neuron in the next layer and tend to capture more global features. 
CNNs are therefore more efficient and effective for processing image data.
One could also note that if we apply a filter that is the same size 
as the input image then the convolution operation becomes equivalent to a 
fully connected (dense) layer therefore CNNs can be thought of as a "generalization" of FNNs.


\subsection{Motivation for CNNs and ML}

Our current methods for estimating the mass of galaxy clusters 
from X-ray data rely on physical models that assume the intracluster gas is in 
hydrostatic equilibrium. While these analytical approaches 
(such as the $\beta$-model) have been very successful, they are 
based on simplifying assumptions, such as spherical symmetry and --
hydrostatic equilibrium, which may not fully capture the complexities of 
real clusters—especially when clusters are disturbed. As a result, 
such models may be limited in accuracy or 
applicability for certain cluster populations.

Machine learning (ML), and Convolutional Neural Networks (CNNs) 
in particular, offer a complementary approach. CNNs are designed 
to extract subtle spatial and spectral patterns from images,
making them ideally suited for X-ray observations.
By training on  datasets of simulated clusters, CNNs can 
learn to recognize features correlated 
with cluster mass that may not be accessible to traditional 
analytical methods. Furthermore, ML approaches can generalize 
to situations where the underlying 
physical assumptions are only approximately valid, or where 
the data is too complex for simple analytic models.
